\section{Conclusion}
\label{sec:conclusion}

This study used mutation testing to examine whether validation workflows
already present in machine-learning research repositories detect controlled
changes in reproducibility-relevant experimental choices. Using MLReproMutate,
we evaluated mutations affecting random seeds, dependency constraints, data
splitting, and cross-validation configuration under an outcome-blind protocol
with frozen repository revisions, mutation candidates, and validation
workflows.

Of the 39 frozen repository--operator cases, 13 were evaluable during primary
execution. Bounded restoration increased the combined evaluable set to 24
cases. After excluding one confirmed-equivalent mutation, 23 confirmed
non-equivalent mutations entered the meaningful detection denominator. The
selected validation workflows detected 2 of these 23 mutations and did not
detect 21.

These results do not imply that the corresponding repositories are
irreproducible. Rather, they show that, for the evaluated cases in this sample,
the selected existing validation workflow did not detect the particular
reproducibility-relevant mutation in 21 of the 23 confirmed non-equivalent
evaluations.

The study also highlights executability as a practical prerequisite for
retrospective mutation analysis. Primary execution left 26 of the 39 frozen
cases non-evaluable; bounded restoration reduced this number to 15 while
preserving the primary execution record.

Our descriptive analysis did not support attributing mutation detection to
oracle explicitness. Both detected B02 mutations occurred in
\texttt{dependency-pin} cases with \texttt{completion-only} workflows, while
small and uneven category sizes, operator composition, and differential
evaluability prevent stronger inference.

More broadly, reproducibility-oriented mutation testing provides a way to ask
a concrete validation question: if an experimentally important choice changes,
does an existing repository safeguard notice? MLReproMutate provides an
auditable mechanism for posing this question through controlled mutations,
baseline-first execution, and explicit handling of semantic equivalence.
